
\documentclass{article}
\usepackage{colortbl}
\usepackage{makecell}
\usepackage{multirow}
\usepackage{supertabular}

\begin{document}

\newcounter{utterance}

\centering \large Interaction Transcript for game `referencegame', experiment `random\_grids', episode 9 with kimi{-}k2.5{-}without{-}reasoning{-}t1.0.
\vspace{24pt}

{ \footnotesize  \setcounter{utterance}{1}
\setlength{\tabcolsep}{0pt}
\begin{supertabular}{c@{$\;$}|p{.15\linewidth}@{}p{.15\linewidth}p{.15\linewidth}p{.15\linewidth}p{.15\linewidth}p{.15\linewidth}}
    \# & $\;$A & \multicolumn{4}{c}{Game Master} & $\;\:$B\\
    \hline

    \theutterance \stepcounter{utterance}  
    & & \multicolumn{4}{p{0.6\linewidth}}{
        \cellcolor[rgb]{0.9,0.9,0.9}{
            \makecell[{{p{\linewidth}}}]{
                \texttt{\tiny{[P1$\langle$GM]}}
                \texttt{You are given three grids, where each of them is 5 by 5 in size.} \\
\texttt{Grids have empty cells marked with "▢" and filled cells marked with "X".} \\
\texttt{Your task is to generate a referring expression that best describes the target grid while distinguishing it from the two other distractor grids.} \\
\texttt{The first grid is the target grid, and the following two grids are the distractors.} \\
\\ 
\texttt{Target grid:} \\
\\ 
\texttt{X □ X X □} \\
\texttt{□ X □ X X} \\
\texttt{X X □ □ X} \\
\texttt{□ X X □ X} \\
\texttt{X □ X □ □} \\
\\ 
\texttt{Distractor grid 1:} \\
\\ 
\texttt{□ □ X X □} \\
\texttt{X □ □ X X} \\
\texttt{□ X X □ X} \\
\texttt{X X □ □ X} \\
\texttt{□ X X □ □} \\
\\ 
\texttt{Distractor grid 2:} \\
\\ 
\texttt{X X □ X □} \\
\texttt{□ X X □ X} \\
\texttt{X □ □ X □} \\
\texttt{□ X □ □ X} \\
\texttt{X □ X □ X} \\
\\ 
\texttt{Instruction: Describe the target grid.} \\
\texttt{Generate the referring expression starting with the tag "Expression: " for the given target grid. Omit any other text.} \\
            }
        }
    }
    & \\ \\

    \theutterance \stepcounter{utterance}  
    & \multicolumn{4}{p{0.6\linewidth}}{
        \cellcolor[rgb]{0.9,0.9,0.9}{
            \makecell[{{p{\linewidth}}}]{
                \texttt{\tiny{[P1$\rangle$GM]}}
                \texttt{Expression: The grid where the top{-}left corner has an X, the top{-}right corner is empty, and there is a diagonal line of X's running from the top{-}left area down through the center, with X's at positions (1,1), (2,2), (3,3), (4,4), (5,5) being mostly filled except for (3,3) which is empty, and notably has X's in the first column at rows 1, 3, and 5.} \\
            }
        }
    }
    & & \\ \\

    \theutterance \stepcounter{utterance}  
    & & & \multicolumn{2}{p{0.3\linewidth}}{
        \cellcolor[rgb]{0.9,0.9,0.9}{
            \makecell[{{p{\linewidth}}}]{
                \texttt{\tiny{[GM$|$GM]}}
                \texttt{Expression: The grid where the top{-}left corner has an X, the top{-}right corner is empty, and there is a diagonal line of X's running from the top{-}left area down through the center, with X's at positions (1,1), (2,2), (3,3), (4,4), (5,5) being mostly filled except for (3,3) which is empty, and notably has X's in the first column at rows 1, 3, and 5.} \\
            }
        }
    }
    & & \\ \\

    \theutterance \stepcounter{utterance}  
    & & \multicolumn{4}{p{0.6\linewidth}}{
        \cellcolor[rgb]{0.9,0.9,0.9}{
            \makecell[{{p{\linewidth}}}]{
                \texttt{\tiny{[GM$\rangle$P2]}}
                \texttt{You are given three grids, where each of them is 5 by 5 in size.} \\
\texttt{Grids have empty cells marked with "▢" and filled cells marked with "X".} \\
\texttt{You are also given a referring expression that describes one of the given grids.} \\
\texttt{Your task is to select the grid that matches the given referring expression.} \\
\texttt{Generate only the number (in text) of the grid that matches the given expression by selecting first, second, or third.} \\
\\ 
\texttt{First:} \\
\\ 
\texttt{X □ X X □} \\
\texttt{□ X □ X X} \\
\texttt{X X □ □ X} \\
\texttt{□ X X □ X} \\
\texttt{X □ X □ □} \\
\\ 
\texttt{Second:} \\
\\ 
\texttt{□ □ X X □} \\
\texttt{X □ □ X X} \\
\texttt{□ X X □ X} \\
\texttt{X X □ □ X} \\
\texttt{□ X X □ □} \\
\\ 
\texttt{Third:} \\
\\ 
\texttt{X X □ X □} \\
\texttt{□ X X □ X} \\
\texttt{X □ □ X □} \\
\texttt{□ X □ □ X} \\
\texttt{X □ X □ X} \\
\\ 
\texttt{Expression: The grid where the top{-}left corner has an X, the top{-}right corner is empty, and there is a diagonal line of X's running from the top{-}left area down through the center, with X's at positions (1,1), (2,2), (3,3), (4,4), (5,5) being mostly filled except for (3,3) which is empty, and notably has X's in the first column at rows 1, 3, and 5.} \\
\texttt{Question: Which grid does the expression refer to?} \\
\texttt{Start with the tag "Answer: ", followed by your selection. Omit any other text.} \\
            }
        }
    }
    & \\ \\

    \theutterance \stepcounter{utterance}  
    & & & \multicolumn{4}{p{0.6\linewidth}}{
        \cellcolor[rgb]{0.9,0.9,0.9}{
            \makecell[{{p{\linewidth}}}]{
                \texttt{\tiny{[GM$\langle$P2]}}
                \texttt{I need to analyze each grid based on the given criteria.} \\
\\ 
\texttt{Let me check the conditions:} \\
\texttt{1. Top{-}left corner has an X (position 1,1)} \\
\texttt{2. Top{-}right corner is empty (position 1,5)} \\
\texttt{3. Diagonal from top{-}left to bottom{-}right with X's at (1,1), (2,2), (3,3), (4,4), (5,5) {-} but (3,3) should be empty} \\
\texttt{4. X's in first column at rows 1, 3, and 5} \\
\\ 
\texttt{Let me check each grid:} \\
\\ 
\texttt{**First Grid:**} \\
\texttt{```} \\
\texttt{X □ X X □} \\
\texttt{□ X □ X X} \\
\texttt{X X □ □ X} \\
\texttt{□ X X □ X} \\
\texttt{X □ X □ □} \\
\texttt{```} \\
\texttt{{-} (1,1) = X ✓} \\
\texttt{{-} (1,5) = □ ✓} \\
\texttt{{-} Diagonal: (1,1)=X, (2,2)=X, (3,3)=□, (4,4)=□, (5,5)=□ {-} but we need (4,4) and (5,5) to be X, and they aren't ✗} \\
\\ 
\texttt{**Second Grid:**} \\
\texttt{```} \\
\texttt{□ □ X X □} \\
\texttt{X □ □ X X} \\
\texttt{□ X X □ X} \\
\texttt{X X □ □ X} \\
\texttt{□ X X □ □} \\
\texttt{```} \\
\texttt{{-} (1,1) = □ ✗} \\
\\ 
\texttt{**Third Grid:**} \\
\texttt{```} \\
\texttt{X X □ X □} \\
\texttt{□ X X □ X} \\
\texttt{X □ □ X □} \\
\texttt{□ X □ □ X} \\
\texttt{X □ X □ X} \\
\texttt{```} \\
\texttt{{-} (1,1) = X ✓} \\
\texttt{{-} (1,5) = □ ✓} \\
\texttt{{-} Diagonal: (1,1)=X, (2,2)=X, (3,3)=□, (4,4)=□, (5,5)=X {-} but we need (4,4) to be X, and it isn't ✗} \\
\\ 
\texttt{Wait, let me re{-}read. The expression says "X's at positions (1,1), (2,2), (3,3), (4,4), (5,5) being mostly filled except for (3,3) which is empty". This means (1,1), (2,2), (4,4), (5,5) should be X, and (3,3) should be empty.} \\
\\ 
\texttt{Let me recheck Third Grid:} \\
\texttt{{-} (1,1)=X, (2,2)=X, (3,3)=□, (4,4)=□, (5,5)=X} \\
\\ 
\texttt{(4,4) is □, not X. So this doesn't match either.} \\
\\ 
\texttt{Wait, let me recheck First Grid more carefully:} \\
\texttt{{-} (1,1)=X, (2,2)=X, (3,3)=□, (4,4)=□, (5,5)=□} \\
\\ 
\texttt{That doesn't match.} \\
\\ 
\texttt{Let me re{-}read the expression again. It says "diagonal line of X's running from the top{-}left area down through the center" {-} maybe it's not the main diagonal?} \\
\\ 
\texttt{Actually, re{-}reading: "X's at positions (1,1), (2,2), (3,3), (4,4), (5,5) being mostly filled except for (3,3) which is empty"} \\
\\ 
\texttt{For Third Grid: positions (1,1)=X, (2,2)=X, (3,3)=□, (4,4)=□, (5,5)=X {-} that's 3 X's and 2 empty.} \\
\\ 
\texttt{For First Grid: (1,1)=X, (2,2)=X, (3,3)=□, (4,4)=□, (5,5)=□ {-} that's 2 X's and 3 empty.} \\
\\ 
\texttt{Hmm, but wait {-} let me check first column for Third Grid: (1,1)=X, (2,1)=□, (3,1)=X, (4,1)=□, (5,1)=X {-} that's X's at rows 1, 3, and 5! ✓} \\
\\ 
\texttt{And for First Grid first column: X, □, X, □, X {-} also rows 1, 3, 5! ✓} \\
\\ 
\texttt{Going back to the diagonal description {-} "mostly filled except for (3,3)". Third grid has (4,4) empty too.} \\
\\ 
\texttt{But wait, re{-}reading more carefully {-} maybe I misread the third grid. Let me check again:} \\
\texttt{Row 4: □ X □ □ X {-} so (4,4) is indeed □.} \\
\\ 
\texttt{Actually, re{-}reading the expression once more: it says the diagonal has X's at those positions "being mostly filled except for (3,3)". This suggests (4,4) and (5,5) should be filled.} \\
\\ 
\texttt{Hmm, but neither grid perfectly matches. Let me check if I misread something...} \\
\\ 
\texttt{Actually, looking at Third Grid row 5: X □ X □ X {-} so (5,5)=X ✓} \\
\\ 
\texttt{The issue is (4,4). But maybe the description is approximate? Third Grid matches best on the first column requirement and has the most diagonal X's with only (3,3) and (4,4) empty.} \\
\\ 
\texttt{Wait {-} I need to re{-}check. The expression says "except for (3,3) which is empty" {-} implying (4,4) and (5,5) should be X. But in Third Grid, (4,4)=□.} \\
\\ 
\texttt{Let me re{-}check First Grid row 4: □ X X □ X {-} (4,3)=X, (4,4)=□, (4,5)=X} \\
\\ 
\texttt{Hmm, neither has (4,4)=X. } \\
\\ 
\texttt{But Third Grid has (5,5)=X and First Grid has (5,5)=□. The expression mentions (5,5) should be filled. So Third Grid is closer.} \\
\\ 
\texttt{Also checking: Third Grid has first column X at 1,3,5 {-} matches perfectly.} \\
\\ 
\texttt{Answer: third} \\
            }
        }
    }
     \\ \\

    \theutterance \stepcounter{utterance}  
    & & & \multicolumn{2}{p{0.3\linewidth}}{
        \cellcolor[rgb]{0.9,0.9,0.9}{
            \makecell[{{p{\linewidth}}}]{
                \texttt{\tiny{[GM$|$GM]}}
                \texttt{Invalid generated expression} \\
            }
        }
    }
    & & \\ \\

\end{supertabular}
}

\end{document}
